\makeatletter
\def\input@path{{../styles/}{../../styles/}{../../../styles/}{../}{../../}{../../../}}
\makeatother
\documentclass{ee102_notes}
\input{macros.tex}
\input{preamble.tex}
\renewcommand{\releasedate}{September 29, 2025}

\addtolength{\oddsidemargin}{-0.15in}
\addtolength{\evensidemargin}{-0.15in}
\addtolength{\textwidth}{0.30in}
\addtolength{\topmargin}{-0.10in}
\addtolength{\textheight}{0.20in}

\newcommand{\Eblank}{\rule{3cm}{0.4pt}}
\newcommand{\Rankblank}{\rule{3cm}{0.4pt}}

\newcommand{\uof}[1]{u\!\left[#1\right]} % discrete-time unit step

\begin{document}
% \section{In-class activity on convolution}
% ---- knobs you can tweak ----
\def\Alpha{0.5}   % = 1/2
\def\HLen{5}      % h[n] nonzero for n=0..5
\def\Xmin{-2} \def\Xmax{5}    % x[k] axis range
\def\Kmin{-6} \def\Kmax{10}   % h[n-k] and product axis range
\def\xscale{0.30cm} \def\yscale{1.15cm}
\def\panelgap{0.35cm}

% ---- helper: empty axes with labels ----
\newcommand{\EmptyAxes}[4]{%
% #1 xmin, #2 xmax, #3 y-label, #4 x-label (usually k)
\begin{tikzpicture}[x=\xscale,y=\yscale]
  \draw[->] (#1-0.6,0) -- (#2+0.6,0) node[right] {#4};
  \draw[->] (0,-0.05) -- (0,1.05) node[left] {#3};
  \foreach \t in {#1,...,#2} \node[below=2pt,font=\fontsize{5}{5}\selectfont] at (\t,0) {\t};
\end{tikzpicture}%
}

% ---- (A) BLANK active-learning panel for a given n ----
\newcommand{\BlankConvPanel}[1]{%
\begin{minipage}{\linewidth}
\noindent\textbf{Hold $n=#1$ fixed}\\[-2pt]
\EmptyAxes{\Xmin}{\Xmax}{$x[k]$}{$k$}\hspace{\panelgap}%
\EmptyAxes{\Kmin}{\Kmax}{$h[#1-k]$}{$k$}\hspace{\panelgap}%
\EmptyAxes{\Kmin}{\Kmax}{$p_{#1}[k]$}{$k$}\\[-2pt]
{\footnotesize Overlap indices: $k=\rule{2.5cm}{0.4pt}$\hfill
$y[#1]=\sum_k p_{#1}[k]=\rule{2.5cm}{0.4pt}$}\par\vspace{1pt}
\end{minipage}\par
}

% ---- (B) FILLED panel for n=0 with answer + short solution ----
\newcommand{\FilledPanelNZero}{%
\noindent\textbf{Worked row: hold $n=0$ fixed}\\[-2pt]
% x[k]
\begin{tikzpicture}[x=0.65*\xscale,y=0.65*\yscale]
  \draw[->] (\Xmin-0.6,0) -- (\Xmax+0.6,0) node[right] {$k$};
  \draw[->] (0,-0.05) -- (0,1.15) node[left] {$x[k]$};
  \foreach \t in {\Xmin,...,\Xmax} \node[below=2pt,font=\fontsize{5}{5}\selectfont] at (\t,0) {\t};
  % stems: x[0]=x[1]=x[2]=1
  \foreach \k in {0,1,2}{
    \draw[very thick] (\k,0) -- (\k,1);
    \fill (\k,1) circle (1.6pt);
  }
\end{tikzpicture}
\hspace{\panelgap}
% h[n-k] with n=0: nonzero for k in [-5,0], value alpha^{-k}
\begin{tikzpicture}[x=0.65*\xscale,y=0.65*\yscale]
  \draw[->] (\Kmin-0.6,0) -- (\Kmax-4.0,0) node[right] {$k$};
  \draw[->] (0,-0.05) -- (0,1.15) node[left] {$h[0-k]$};
  \foreach \t in {-6,...,2} \node[below=2pt,font=\fontsize{5}{5}\selectfont] at (\t,0) {\t};
  \foreach \k in {-5,-4,...,0}{
    \pgfmathparse{pow(\Alpha,-\k)}
    \let\val\pgfmathresult
    \draw[very thick] (\k,0) -- (\k,\val);
    \fill (\k,\val) circle (1.6pt);
  }
\end{tikzpicture}
\hspace{\panelgap}%
% Pointwise product x[k]h[-k], nonzero only at k=0
\begin{tikzpicture}[x=0.65*\xscale,y=0.65*\yscale]
  \draw[->] (\Xmin-0.6,0) -- (\Xmax+0.6,0) node[right] {\ensuremath{k}};
  \draw[->] (0,-0.05) -- (0,1.15) node[left] {\ensuremath{p_{0}[k]}};
  \foreach \t in {\Xmin,...,\Xmax} \node[below=2pt,font=\fontsize{5}{5}\selectfont] at (\t,0) {\t};
  \draw[very thick] (0,0) -- (0,1);
  \fill (0,1) circle (1.8pt);
\end{tikzpicture}
\par\smallskip
\noindent{\footnotesize Overlap index: \(k=0\). Hence
\(y[0]=\sum_k p_{0}[k]=x[0]h[0]+x[1]h[-1]+x[2]h[-2]=1\).}\par\vspace{1pt}
}

% ---- final output axis ----
\newcommand{\OutputAxes}{%
\begin{tikzpicture}[x=1.05cm,y=0.65cm]
  \draw[->] (-1.6,0) -- (8.7,0) node[right] {$n$};
  \draw[->] (0,-0.08) -- (0,2.2) node[above right] {$y[n]$};
  \foreach \t in {-1,...,8} {
    \draw (\t,0.035) -- (\t,-0.035);
    \node[below=2pt,font=\scriptsize] at (\t,0) {\t};
  }
  \foreach \a in {1,2} {
    \draw (0.04,\a) -- (-0.04,\a);
    \node[left=2pt,font=\scriptsize] at (0,\a) {\a};
  }
\end{tikzpicture}%
}

% ---- (C) Generate the 10 requested panels ----
% Tip: Put five per page using manual page breaks or minipages.
% Page 1 (example):
\thispagestyle{fancy}
% add header: EE 102 activity 
\lhead{NAME: \underline{\hspace{5cm}}}
\rhead{EE 102: In-class activity}
\section*{Visualize convolution}
For $x[n]=u[n]-u[n-3]$ and
\[
h[n]=
\begin{cases}
\left(\tfrac12\right)^n, & 0\le n\le 5,\\
0, & \text{otherwise,}
\end{cases}
\]
For each fixed $n$, define $p_{n}[k]=x[k]h[n-k]$. The index $k$ runs along every graph in a row;
$n$ stays fixed while completing that row.
\begin{multicols}{2}
\begin{enumerate}[leftmargin=1.7em,itemsep=0pt,topsep=1pt]
  \item Draw $x[k]$.
  \item Flip and shift $h[k]$ to draw $h[n-k]$.
  \item Mark the overlap and draw $p_{n}[k]$.
  \item Add the product samples: $y[n]=\sum_k p_{n}[k]$.
\end{enumerate}
\end{multicols}
Use $0\le k\le2$ for the support of $x[k]$ and $n-5\le k\le n$ for the
support of $h[n-k]$.
\medskip
\noindent\textbf{At $n=-1$:} No overlap occurs, so
$x[k]h[-1-k]=0$ for all $k$, giving $y[-1]=0$.
\medskip
\hrule
\FilledPanelNZero
\hrule
\BlankConvPanel{1}
\hrule
\BlankConvPanel{2}
\hrule
\newpage
\BlankConvPanel{3}
\hrule
\BlankConvPanel{4}
\hrule
\BlankConvPanel{5}
\hrule
\BlankConvPanel{6}
\hrule
\BlankConvPanel{7}
\medskip
\begin{minipage}{\linewidth}
\noindent\textbf{Assemble the output}\\[-2pt]

Plot the samples from $n=-1$ through $n=8$, including the zero-valued
samples. The nonzero support is $\rule{1.2cm}{0.4pt}\le n\le\rule{1.2cm}{0.4pt}$.
\begin{center}
\OutputAxes
\end{center}
\end{minipage}
\end{document}
